\section{AHSME 1950}
        \begin{enumerate}
        	\item (\textit{AHSME 1950}) If $64$ is divided into three parts proportional to $2$, $4$, and $6$, the smallest part is:


 $\textbf{(A)}\ 5\frac{1}{3}\qquad\textbf{(B)}\ 11\qquad\textbf{(C)}\ 10\frac{2}{3}\qquad\textbf{(D)}\ 5\qquad\textbf{(E)}\ \text{None of these answers}$


 

	\item (\textit{AHSME 1950}) Let $R=gS-4$. When $S=8$, $R=16$. When $S=10$, $R$ is equal to:


 $\textbf{(A)}\ 11\qquad\textbf{(B)}\ 14\qquad\textbf{(C)}\ 20\qquad\textbf{(D)}\ 21\qquad\textbf{(E)}\ \text{None of these}$


 

	\item (\textit{AHSME 1950}) The sum of the roots of the equation $4x^{2}+5-8x=0$ is equal to:


 $\textbf{(A)}\ 8\qquad\textbf{(B)}\ -5\qquad\textbf{(C)}\ -\frac{5}{4}\qquad\textbf{(D)}\ -2\qquad\textbf{(E)}\ \text{None of these}$


 

	\item (\textit{AHSME 1950}) Reduced to lowest terms, $\frac{a^{2}-b^{2}}{ab} - \frac{ab-b^{2}}{ab-a^{2}}$ is equal to:


 $\textbf{(A)}\ \frac{a}{b}\qquad\textbf{(B)}\ \frac{a^{2}-2b^{2}}{ab}\qquad\textbf{(C)}\ a^{2}\qquad\textbf{(D)}\ a-2b\qquad\textbf{(E)}\ \text{None of these}$


 

	\item (\textit{AHSME 1950}) If five geometric means are inserted between $8$ and $5832$, the fifth term in the geometric series is:


 $\textbf{(A)}\ 648\qquad\textbf{(B)}\ 832\qquad\textbf{(C)}\ 1168\qquad\textbf{(D)}\ 1944\qquad\textbf{(E)}\ \text{None of these}$


 

	\item (\textit{AHSME 1950}) The values of $y$ which will satisfy the equations 
 \[\begin{array}{rcl} 2x^{2}+6x+5y+1&=&0\\ 2x+y+3&=&0 \end{array}\] may be found by solving:


 $\textbf{(A)}\ y^{2}+14y-7=0\qquad\textbf{(B)}\ y^{2}+8y+1=0\qquad\textbf{(C)}\ y^{2}+10y-7=0\qquad\\ \textbf{(D)}\ y^{2}+y-12=0\qquad \textbf{(E)}\ \text{None of these equations}$


 

	\item (\textit{AHSME 1950}) If the digit $1$ is placed after a two digit number whose tens' digit is $t$, and units' digit is $u$, the new number is:


 $\textbf{(A)}\ 10t+u+1\qquad\textbf{(B)}\ 100t+10u+1\qquad\textbf{(C)}\ 1000t+10u+1\qquad\textbf{(D)}\ t+u+1\qquad\\ \textbf{(E)}\ \text{None of these answers}$


 

	\item (\textit{AHSME 1950}) If the radius of a circle is increased $100\%$, the area is increased:


 $\textbf{(A)}\ 100\%\qquad\textbf{(B)}\ 200\%\qquad\textbf{(C)}\ 300\%\qquad\textbf{(D)}\ 400\%\qquad\textbf{(E)}\ \text{By none of these}$


 

	\item (\textit{AHSME 1950}) The largest area of a triangle that can be inscribed in a semi-circle whose radius is $r$ is:


 $\textbf{(A)}\ r^{2}\qquad\textbf{(B)}\ r^{3}\qquad\textbf{(C)}\ 2r^{2}\qquad\textbf{(D)}\ 2r^{3}\qquad\textbf{(E)}\ \frac{1}{2}r^{2}$


 

	\item (\textit{AHSME 1950}) After rationalizing the numerator of $\frac{\sqrt{3}-\sqrt{2}}{\sqrt{3}}$, the denominator in simplest form is:


 $\textbf{(A)}\ \sqrt{3}(\sqrt{3}+\sqrt{2})\qquad\textbf{(B)}\ \sqrt{3}(\sqrt{3}-\sqrt{2})\qquad\textbf{(C)}\ 3-\sqrt{3}\sqrt{2}\qquad\\ \textbf{(D)}\ 3+\sqrt6\qquad\textbf{(E)}\ \text{None of these answers}$


 

	\item (\textit{AHSME 1950}) If in the formula $C =\frac{en}{R+nr}$, $n$ is increased while $e$, $R$ and $r$ are kept constant, then $C$:


 $\textbf{(A)}\ \text{decreases}\qquad\textbf{(B)}\ \text{increases}\qquad\textbf{(C)}\ \text{remains constant}\qquad\textbf{(D)}\ \text{increases and then decreases}\qquad\\ \textbf{(E)}\ \text{decreases and then increases}$


 

	\item (\textit{AHSME 1950}) As the number of sides of a polygon increases from $3$ to $n$, the sum of the exterior angles formed by extending each side in succession:


 $\textbf{(A)}\ \text{increases}\qquad\textbf{(B)}\ \text{decreases}\qquad\textbf{(C)}\ \text{remains constant}\qquad\textbf{(D)}\ \text{cannot be predicted}\qquad\\ \textbf{(E)}\ \text{becomes }(n-3)\text{ straight angles}$


 

	\item (\textit{AHSME 1950}) The roots of $(x^{2}-3x+2)(x)(x-4)=0$ are:


 $\textbf{(A)}\ 4\qquad\textbf{(B)}\ 0\text{ and }4\qquad\textbf{(C)}\ 1\text{ and }2\qquad\textbf{(D)}\ 0,1,2\text{ and }4\qquad\textbf{(E)}\ 1,2\text{ and }4$


 

	\item (\textit{AHSME 1950}) For the simultaneous equations 
 \[2x-3y=8\]
 \[6y-4x=9\]


 $\textbf{(A)}\ x=4,y=0\qquad\textbf{(B)}\ x=0,y=\frac{3}{2}\qquad\textbf{(C)}\ x=0,y=0\qquad\\ \textbf{(D)}\ \text{There is no solution}\qquad\textbf{(E)}\ \text{There are an infinite number of solutions}$


 

	\item (\textit{AHSME 1950}) The real factors of $x^2+4$ are:


 $\textbf{(A)}\ (x^{2}+2)(x^{2}+2)\qquad\textbf{(B)}\ (x^{2}+2)(x^{2}-2)\qquad\textbf{(C)}\ x^{2}(x^{2}+4)\qquad\\ \textbf{(D)}\ (x^{2}-2x+2)(x^{2}+2x+2)\qquad\textbf{(E)}\ \text{Non-existent}$


 

	\item (\textit{AHSME 1950}) The number of terms in the expansion of $[(a+3b)^{2}(a-3b)^{2}]^{2}$ when simplified is:


 $\textbf{(A)}\ 4\qquad\textbf{(B)}\ 5\qquad\textbf{(C)}\ 6\qquad\textbf{(D)}\ 7\qquad\textbf{(E)}\ 8$


 

	\item (\textit{AHSME 1950}) The formula which expresses the relationship between $x$ and $y$ as shown in the accompanying table is:


 
 \[\begin{tabular}[t]{|c|c|c|c|c|c|}\hline x&0&1&2&3&4\\\hline y&100&90&70&40&0\\\hline\end{tabular}\]


 $\textbf{(A)}\ y=100-10x\qquad\textbf{(B)}\ y=100-5x^{2}\qquad\textbf{(C)}\ y=100-5x-5x^{2}\qquad\\ \textbf{(D)}\ y=20-x-x^{2}\qquad\textbf{(E)}\ \text{None of these}$


 

	\item (\textit{AHSME 1950}) Of the following 


 (1) $a(x-y)=ax-ay$


 (2) $a^{x-y}=a^x-a^y$


 (3) $\log (x-y)=\log x-\log y$


 (4) $\frac{\log x}{\log y}=\log{x}-\log{y}$


 (5) $a(xy)=ax \cdot ay$


 $\textbf{(A)}\ \text{Only 1 and 4 are true}\qquad\\\textbf{(B)}\ \text{Only 1 and 5 are true}\qquad\\\textbf{(C)}\ \text{Only 1 and 3 are true}\qquad\\\textbf{(D)}\ \text{Only 1 and 2 are true}\qquad\\\textbf{(E)}\ \text{Only 1 is true}$


 

	\item (\textit{AHSME 1950}) If $m$ men can do a job in $d$ days, then $m+r$ men can do the job in:


 $\textbf{(A)}\ d+r\text{ days}\qquad\textbf{(B)}\ d-r\text{ days}\qquad\textbf{(C)}\ \frac{md}{m+r}\text{ days}\qquad\textbf{(D)}\ \frac{d}{m+r}\text{ days}\qquad\textbf{(E)}\ \text{None of these}$


 

	\item (\textit{AHSME 1950}) When $x^{13}+1$ is divided by $x-1$, the remainder is:


 $\textbf{(A)}\ 1\qquad\textbf{(B)}\ -1\qquad\textbf{(C)}\ 0\qquad\textbf{(D)}\ 2\qquad\textbf{(E)}\ \text{None of these answers}$


 

	\item (\textit{AHSME 1950}) The volume of a rectangular solid each of whose side, front, and bottom faces are $12\text{ in}^{2}$, $8\text{ in}^{2}$, and $6\text{ in}^{2}$ respectively is:


 $\textbf{(A)}\ 576\text{ in}^{3}\qquad\textbf{(B)}\ 24\text{ in}^{3}\qquad\textbf{(C)}\ 9\text{ in}^{3}\qquad\textbf{(D)}\ 104\text{ in}^{3}\qquad\textbf{(E)}\ \text{None of these}$


 

	\item (\textit{AHSME 1950}) Successive discounts of $10\%$ and $20\%$ are equivalent to a single discount of:


 $\textbf{(A)}\ 30\%\qquad\textbf{(B)}\ 15\%\qquad\textbf{(C)}\ 72\%\qquad\textbf{(D)}\ 28\%\qquad\textbf{(E)}\ \text{None of these}$


 

	\item (\textit{AHSME 1950}) A man buys a house for $10000$ dollars and rents it. He puts $12\frac{1}{2}\%$ of each month's rent aside for repairs and upkeep; pays $325$ dollars a year taxes and realizes $5\frac{1}{2}\%$ on his investment. The monthly rent is:


 $\textbf{(A)}\ \$64.82 \qquad\textbf{(B)}\ \$83.33 \qquad\textbf{(C)}\ \$72.08\qquad\textbf{(D)}\ \$45.83 \qquad\textbf{(E)}\ \$177.08$


 

	\item (\textit{AHSME 1950}) The equation $x+\sqrt{x-2}=4$ has:


 $\textbf{(A)}\ \text{2 real roots}\qquad\textbf{(B)}\ \text{1 real and 1 imaginary root}\qquad\textbf{(C)}\ \text{2 imaginary roots}\qquad\textbf{(D)}\ \text{No roots}\qquad\textbf{(E)}\ \text{1 real root}$


 

	\item (\textit{AHSME 1950}) The value of $\log_{5}\frac{(125)(625)}{25}$ is equal to:


 $\textbf{(A)}\ 725\qquad\textbf{(B)}\ 6\qquad\textbf{(C)}\ 3125\qquad\textbf{(D)}\ 5\qquad\textbf{(E)}\ \text{None of these}$


 

	\item (\textit{AHSME 1950}) if $\log_{10}{m}= b-\log_{10}{n}$, then $m=$


 $\textbf{(A)}\ \frac{b}{n}\qquad\textbf{(B)}\ bn\qquad\textbf{(C)}\ 10^{b}n\qquad\textbf{(D)}\ b-10^{n}\qquad\textbf{(E)}\ \frac{10^{b}}{n}$


 

	\item (\textit{AHSME 1950}) A car travels $120$ miles from $A$ to $B$ at $30$ miles per hour but returns the same distance at $40$ miles per hour. The average speed for the round trip is closest to:


 $\textbf{(A)}\ 33\text{ mph}\qquad\textbf{(B)}\ 34\text{ mph}\qquad\textbf{(C)}\ 35\text{ mph}\qquad\textbf{(D)}\ 36\text{ mph}\qquad\textbf{(E)}\ 37\text{ mph}$


 

	\item (\textit{AHSME 1950}) Two boys $A$ and $B$ start at the same time to ride from Port Jervis to Poughkeepsie, $60$ miles away. $A$ travels $4$ miles an hour slower than $B$.  $B$ reaches Poughkeepsie and at once turns back meeting $A$ $12$ miles from Poughkeepsie. The rate of $A$ was:


 $\textbf{(A)}\ 4\text{ mph}\qquad\textbf{(B)}\ 8\text{ mph}\qquad\textbf{(C)}\ 12\text{ mph}\qquad\textbf{(D)}\ 16\text{ mph}\qquad\textbf{(E)}\ 20\text{ mph}$


 

	\item (\textit{AHSME 1950}) A manufacturer built a machine which will address $500$ envelopes in $8$ minutes. He wishes to build another machine so that when both are operating together they will address $500$ envelopes in $2$ minutes. The equation used to find how many minutes $x$ it would require the second machine to address $500$ envelopes alone is:


 $\textbf{(A)}\ 8-x=2\qquad\textbf{(B)}\ \frac{1}{8}+\frac{1}{x}=\frac{1}{2}\qquad\textbf{(C)}\ \frac{500}{8}+\frac{500}{x}=500\qquad\textbf{(D)}\ \frac{x}{2}+\frac{x}{8}=1\qquad\\ \textbf{(E)}\ \text{none of these answers}$


 

	\item (\textit{AHSME 1950}) From a group of boys and girls, $15$ girls leave. There are then left two boys for each girl. After this $45$ boys leave. There are then $5$ girls for each boy. The number of girls in the beginning was:


 $\textbf{(A)}\ 40\qquad\textbf{(B)}\ 43\qquad\textbf{(C)}\ 29\qquad\textbf{(D)}\ 50\qquad\textbf{(E)}\ \text{None of these}$


 

	\item (\textit{AHSME 1950}) John ordered $4$ pairs of black socks and some additional pairs of blue socks. The price of the black socks per pair was twice that of the blue. When the order was filled, it was found that the number of pairs of the two colors had been interchanged. This increased the bill by $50\%$. The ratio of the number of pairs of black socks to the number of pairs of blue socks in the original order was:


 $\textbf{(A)}\ 4:1\qquad\textbf{(B)}\ 2:1\qquad\textbf{(C)}\ 1:4\qquad\textbf{(D)}\ 1:2\qquad\textbf{(E)}\ 1:8$


 

	\item (\textit{AHSME 1950}) A $25$ foot ladder is placed against a vertical wall of a building. The foot of the ladder is $7$ feet from the base of the building. If the top of the ladder slips $4$ feet, then the foot of the ladder will slide:


 $\textbf{(A)}\ 9\text{ ft}\qquad\textbf{(B)}\ 15\text{ ft}\qquad\textbf{(C)}\ 5\text{ ft}\qquad\textbf{(D)}\ 8\text{ ft}\qquad\textbf{(E)}\ 4\text{ ft}$


 

	\item (\textit{AHSME 1950}) The number of circular pipes with an inside diameter of $1$ inch which will carry the same amount of water as a pipe with an inside diameter of $6$ inches is:


 $\textbf{(A)}\ 6\pi\qquad\textbf{(B)}\ 6\qquad\textbf{(C)}\ 12\qquad\textbf{(D)}\ 36\qquad\textbf{(E)}\ 36\pi$


 

	\item (\textit{AHSME 1950}) When the circumference of a toy balloon is increased from $20$ inches to $25$ inches, the radius is increased by:


 $\textbf{(A)}\ 5\text{ in}\qquad\textbf{(B)}\ 2\frac{1}{2}\text{ in}\qquad\textbf{(C)}\ \frac{5}{\pi}\text{ in}\qquad\textbf{(D)}\ \frac{5}{2\pi}\text{ in}\qquad\textbf{(E)}\ \frac{\pi}{5}\text{ in}$


 

	\item (\textit{AHSME 1950}) In triangle $ABC$, $AC=24$ inches, $BC=10$ inches, $AB=26$ inches. The radius of the inscribed circle is:


 $\textbf{(A)}\ 26\text{ in}\qquad\textbf{(B)}\ 4\text{ in}\qquad\textbf{(C)}\ 13\text{ in}\qquad\textbf{(D)}\ 8\text{ in}\qquad\textbf{(E)}\ \text{None of these}$


 

	\item (\textit{AHSME 1950}) A merchant buys goods at $25\%$ off the list price. He desires to mark the goods so that he can give a discount of $20\%$ on the marked price and still clear a profit of $25\%$ on the selling price. What percent of the list price must he mark the goods?


 $\textbf{(A)}\ 125\%\qquad\textbf{(B)}\ 100\%\qquad\textbf{(C)}\ 120\%\qquad\textbf{(D)}\ 80\%\qquad\textbf{(E)}\ 75\%$


 

	\item (\textit{AHSME 1950}) If $y =\log_{a}{x}$, $a>1$, which of the following statements is incorrect?


 $\textbf{(A)}\ \text{If }x=1,y=0\qquad\\ \textbf{(B)}\ \text{If }x=a,y=1\qquad\\ \textbf{(C)}\ \text{If }x=-1,y\text{ is imaginary (complex)}\qquad\\ \textbf{(D)}\ \text{If }0<x<1,y\text{ is always less than 0 and decreases without limit as }x\text{ approaches zero}\qquad\\ \textbf{(E)}\ \text{Only some of the above statements are correct}$


 

	\item (\textit{AHSME 1950}) If the expression $\begin{vmatrix}a & c\\ d & b\end{vmatrix}$ has the value $ab-cd$ for all values of $a, b, c$ and $d$, then the equation $\begin{vmatrix}2x & 1\\ x & x\end{vmatrix}= 3$:


 $\textbf{(A)}\ \text{Is satisfied for only 1 value of }x\qquad\\ \textbf{(B)}\ \text{Is satisified for 2 values of }x\qquad\\ \textbf{(C)}\ \text{Is satisified for no values of }x\qquad\\ \textbf{(D)}\ \text{Is satisfied for an infinite number of values of }x\qquad\\ \textbf{(E)}\ \text{none of these}$


 

	\item (\textit{AHSME 1950}) Given the series $2+1+\frac{1}{2}+\frac{1}{4}+...$ and the following five statements:


 (1) the sum increases without limit


 (2) the sum decreases without limit


 (3) the difference between any term of the sequence and zero can be made less than any positive quantity no matter how small


 (4) the difference between the sum and 4 can be made less than any positive quantity no matter how small


 (5) the sum approaches a limit


 Of these statements, the correct ones are:


 $\textbf{(A)}\ \text{Only }3\text{ and }4\qquad\textbf{(B)}\ \text{Only }5\qquad\textbf{(C)}\ \text{Only }2\text{ and }4\qquad\textbf{(D)}\ \text{Only }2,3\text{ and }4\qquad\textbf{(E)}\ \text{Only }4\text{ and }5$


 

	\item (\textit{AHSME 1950}) The limit of $\frac{x^{2}-1}{x-1}$ as $x$ approaches $1$ as a limit is:


 $\textbf{(A)}\ 0\qquad\textbf{(B)}\ \text{Indeterminate}\qquad\textbf{(C)}\ x-1\qquad\textbf{(D)}\ 2\qquad\textbf{(E)}\ 1$


 

	\item (\textit{AHSME 1950}) The least value of the function $ax^2+bx+c$ ($a>0$) is:


 $\textbf{(A)}\ -\frac{b}{a}\qquad\textbf{(B)}\ -\frac{b}{2a}\qquad\textbf{(C)}\ b^{2}-4ac\qquad\textbf{(D)}\ \frac{4ac-b^{2}}{4a}\qquad\textbf{(E)}\ \text{None of these}$


 

	\item (\textit{AHSME 1950}) The equation $x^{x^{x^{.^{.^.}}}}=2$ is satisfied when $x$ is equal to:


 $\textbf{(A)}\ \infty\qquad\textbf{(B)}\ 2\qquad\textbf{(C)}\ \sqrt[4]{2}\qquad\textbf{(D)}\ \sqrt{2}\qquad\textbf{(E)}\ \text{None of these}$


 

	\item (\textit{AHSME 1950}) The sum to infinity of $\frac{1}{7}+\frac{2}{7^{2}}+\frac{1}{7^{3}}+\frac{2}{7^{4}}+...$ is:


 $\textbf{(A)}\ \frac{1}{5}\qquad\textbf{(B)}\ \frac{1}{24}\qquad\textbf{(C)}\ \frac{5}{48}\qquad\textbf{(D)}\ \frac{1}{16}\qquad\textbf{(E)}\ \text{None of these}$


 

	\item (\textit{AHSME 1950}) The graph of $y=\log x$


 $\textbf{(A)}\ \text{Cuts the }y\text{-axis}\qquad\\ \textbf{(B)}\ \text{Cuts all lines perpendicular to the }x\text{-axis}\qquad\\ \textbf{(C)}\ \text{Cuts the }x\text{-axis}\qquad\\ \textbf{(D)}\ \text{Cuts neither axis}\qquad\\ \textbf{(E)}\ \text{Cuts all circles whose center is at the origin}$


 

	\item (\textit{AHSME 1950}) The number of diagonals that can be drawn in a polygon of $100$ sides is:


 $\textbf{(A)}\ 4850\qquad\textbf{(B)}\ 4950\qquad\textbf{(C)}\ 9900\qquad\textbf{(D)}\ 98\qquad\textbf{(E)}\ 8800$


 

	\item (\textit{AHSME 1950}) In triangle $ABC$, $AB=12$, $AC=7$, and $BC=10$. If sides $AB$ and $AC$ are doubled while $BC$ remains the same, then:


 $\textbf{(A)}\ \text{The area is doubled}\qquad\\ \textbf{(B)}\ \text{The altitude is doubled}\qquad\\ \textbf{(C)}\ \text{The area is four times the original area}\qquad\\ \textbf{(D)}\ \text{The median is unchanged}\qquad\\ \textbf{(E)}\ \text{The area of the triangle is 0}$


 

	\item (\textit{AHSME 1950}) A rectangle inscribed in a triangle has its base coinciding with the base $b$ of the triangle. If the altitude of the triangle is $h$, and the altitude $x$ of the rectangle is half the base of the rectangle, then:


 $\textbf{(A)}\ x=\frac{1}{2}h\qquad\textbf{(B)}\ x=\frac{bh}{b+h}\qquad\textbf{(C)}\ x=\frac{bh}{2h+b}\qquad\textbf{(D)}\ x=\sqrt{\frac{hb}{2}}\qquad\\ \textbf{(E)}\ x=\frac{1}{2}b$


 

	\item (\textit{AHSME 1950}) A point is selected at random inside an equilateral triangle. From this point perpendiculars are dropped to each side. The sum of these perpendiculars is:


 $\textbf{(A)}\ \text{Least when the point is the center of gravity of the triangle}\qquad\\ \textbf{(B)}\ \text{Greater than the altitude of the triangle}\qquad\\ \textbf{(C)}\ \text{Equal to the altitude of the triangle}\qquad\\ \textbf{(D)}\ \text{One-half the sum of the sides of the triangle}\qquad\\ \textbf{(E)}\ \text{Greatest when the point is the center of gravity}$


 

	\item (\textit{AHSME 1950}) A triangle has a fixed base $AB$ that is $2$ inches long. The median from $A$ to side $BC$ is $1 \frac{1}{2}$ inches long and can have any position emanating from $A$. The locus of the vertex $C$ of the triangle is:


 $\textbf{(A)}\ \text{A straight line }AB,1\frac{1}{2}\text{ inches from }A\qquad\\ \textbf{(B)}\ \text{A circle with }A\text{ as center and radius }2\text{ inches}\qquad\\ \textbf{(C)}\ \text{A circle with }A\text{ as center and radius }3\text{ inches}\qquad\\ \textbf{(D)}\ \text{A circle with radius }3\text{ inches and center }4\text{ inches from }B\text{ along }BA\qquad\\ \textbf{(E)}\ \text{An ellipse with }A\text{ as focus}$


 

	\item (\textit{AHSME 1950}) A privateer discovers a merchantman $10$ miles to leeward at $11\text{:}45 \text{ a.m.}$ and with a good breeze bears down upon her at $11$ mph, while the merchantman can only make $8$ mph in his attempt to escape. After a two hour chase, the top sail of the privateer is carried away; she can now make only $17$ miles while the merchantman makes $15$. The privateer will overtake the merchantman at:


 $\textbf{(A)}\ 3\text{:}45\text{ p.m.}\qquad\textbf{(B)}\ 3\text{:}30\text{ p.m.}\qquad\textbf{(C)}\ 5\text{:}00\text{ p.m.}\qquad\textbf{(D)}\ 2\text{:}45\text{ p.m.}\qquad\textbf{(E)}\ 5\text{:}30\text{ p.m.}$


 

\end{enumerate}